\section{Introduction}
\label{sec:intro}

Metaheuristics that learn while they search are a prominent line of work in operations research \citep{CalvetEtal2017,KarimiMamaghanEtal2022}. A typical \emph{learnheuristic} embeds a
reinforcement-learning agent, most often tabular Q-Learning \citep{WatkinsDayan1992,SuttonBarto2018}, as an adaptive operator or strategy selector inside a search framework such as
Iterated Local Search (ILS) \citep{LourencaoEtal2019} or Iterated Greedy (IG) \citep{RuizStutzle2007}, and reports that it beats the search without the agent or published baselines
\citep{KarimiMamaghanEtal2023,MissaouiOSullivan2026}. Such results are usually read as evidence that the \emph{learned policy} is what helps.

Three confounds make that reading fragile. First, the \emph{search structure} that hosts the agent (which operators, which acceptance rule, whether there is a local search) may produce
most of the gain with any selection rule. Second, the \emph{budget}: an agent changes the CPU time per iteration, a local search multiplies it, and comparing methods at an equal number of
iterations can reverse conclusions drawn at equal CPU time. Third, the \emph{tuning}: a parameter validated offline on separate instances is a competitor that most studies omit. Adaptive operator
selection itself is a well-studied topic \citep{FialhoEtal2010,PisingerRopke2007}, but the question ``what does the learned policy add, over what, at what cost?'' is seldom answered with controls
at a matched CPU budget, offline-tuned baselines, certified bounds and hypotheses fixed before the data.

This article answers it with a \emph{budget-fair attribution protocol} (Section~\ref{sec:protocol}) and applies it to three scheduling problems: identical parallel machines
(P$||C_{\max}$), unrelated parallel machines (R$||C_{\max}$) and the permutation flowshop, in which the agent selects the destruction strategy or the destruction size of an IG, or the
operator of an ILS. The evidence is large (twelve synthetic classes, more than $13{,}000$ published identical-machine instances, the standard unrelated-machine benchmark of
\citealt{FanjulPeyroRuiz2010} with third-party reference solutions) and every analysis was specified, with hashes of the code, before the data it uses.

The main findings are the following. (a) The gain of the Q-Learning-guided IG over the ILS comes from the search structure; the learned policy adds $0.01$--$0.05$ percentage points for
identical machines and about $0.1\%$ for unrelated machines, and nothing for the flowshop. (b) The \emph{active ingredient} differs by problem and is revealed only by ablations at equal CPU
time: destruction--reconstruction for identical machines, its combination with a descent for unrelated machines, the local search for the flowshop; a local search that looked $0.15\%$ better
at an equal number of iterations is $0.01\%$ better at equal CPU time. (c) A configuration chosen \emph{offline} on separate instances beats the online agent in all seven environments studied,
against five learners; the headroom of the actions predicts where a policy could help for destruction \emph{strategies} but failed prospectively for the destruction \emph{size}, so headroom is
necessary but not sufficient. (d) Learning-free solvers follow from these findings and are competitive: an IG with a validated destruction size and a start from the best constructive solution
reaches $97.5\%$ of provably optimal solutions on synthetic identical-machine instances, and an IG started from the linear relaxation and polished by a restricted MILP is within $0.25$--$0.30\%$ of the
best known solutions of the standard unrelated-machine benchmark.

The contributions are, in order of strength:
(i) the protocol, with evidence that it changes conclusions (Sections~\ref{sec:protocol} and~\ref{sec:f3f4-ablation}, Table~\ref{tab:fair_vs_naive});
(ii) a controlled decomposition for three problems and two mechanisms, with the active ingredient and the offline-versus-online comparison for each (Sections~\ref{sec:f3f4-ablation},
\ref{sec:unrelated} and~\ref{sec:pfsp}), and a synthesis across seven environments;
(iii) learning-free solvers whose quality is measured against certified bounds and third-party references (Sections~\ref{sec:recipe} and~\ref{sec:unrelated});
(iv) a reproducibility package: code, frozen hashes, a dated research diary with every non-confirmation, and automatic audits that recompute the numbers of the text from the result files.
The article does \emph{not} claim state-of-the-art solvers (exact methods are stronger in both parallel-machine problems), nor that learning cannot help: the agent is tabular, the budgets are those of
the original studies, and richer learners or action spaces are outside the evidence.

Section~\ref{sec:related} reviews related work, Section~\ref{sec:problem} defines the problems, mechanisms and the protocol, Section~\ref{sec:experiments} reports the experiments,
Section~\ref{sec:discussion} discusses their implications and limits and Section~\ref{sec:conclusions} concludes. Further results (comparison with published baselines, robustness checks, granularity and
dispatch, a companion study of ILS) are in the supplementary material.
